% Worked example — self-contained (no external figures), compiles with:
%   latexmk -xelatex slides.tex     (or: xelatex slides.tex, twice)
% Demonstrates the theme + most composition patterns from references/.
\documentclass[aspectratio=169, 11pt, t]{beamer}

\usepackage{beamerthemeAcademicTalk}
\useblue

\usepackage{fontspec}
\usepackage{amsmath, amssymb}
\usepackage{booktabs}
\usetikzlibrary{arrows.meta, positioning, calc}

\setlength{\emergencystretch}{2em}

\title[ERK waves]{A minimal mechanochemical feedback explains ERK waves}
\subtitle{From observation to a testable prediction}
\author[J. Roe]{Jane Roe}
\institute[Example Univ.]{Department of Systems Biology, Example University}
\date{\today}
\setvenue{Biophysics Seminar}

\begin{document}

\begin{frame}[plain]\titlepage\end{frame}

\begin{frame}
  \frametitle{Outline}
  \vskip0.3cm
  {\footnotesize
  \begin{tabbing}
  \hspace{0.4cm}\=\hspace{0.6cm}\=\kill
  \textbf{\color{accentcolor}1}\>\textbf{Question}\,---\,what sustains ERK waves?\\[8pt]
  \textbf{\color{accentcolor}2}\>\textbf{Observation}\,---\,ERK leads cell extension\\[8pt]
  \textbf{\color{accentcolor}3}\>\textbf{Mechanism}\,---\,a two-variable feedback\\[8pt]
  \textbf{\color{accentcolor}4}\>\textbf{Prediction \& test}\,---\,a fixed phase lead\\
  \end{tabbing}}
\end{frame}

% ===== 1 · Question =====
\sectiondivider{1}{Question}

\begin{frame}
  \frametitle{ERK waves coordinate collective migration}
  \vskip0.1cm
  In migrating epithelial monolayers, ERK activity travels as waves whose fronts
  move opposite to cell motion, organizing extension over hundreds of microns.
  \vskip0.2cm
  \keybox{\textbf{Question:} is mechanical deformation alone enough to \alert{sustain}
  and propagate these waves?}
\end{frame}

\begin{frame}
  \frametitle{Existing chemical models predict the wrong phase}
  \vskip0.1cm
  A purely reaction--diffusion account reproduces \emph{a} wave, but places ERK
  activity \alert{behind} strain --- the opposite of what is measured.
  \vskip0.2cm
  That single sign error is the gap this talk closes.
\end{frame}

% ===== 2 · Observation =====
\sectiondivider{2}{Observation}

\begin{frame}
  \frametitle{ERK activation precedes extension by a fixed lag}
  \begin{columns}[T, onlytextwidth]
    \column{0.54\textwidth}
      \vskip0.1cm
      Cross-correlating the ERK and strain channels across fields gives one
      consistent relationship:\par\vskip0.15cm
      $\bullet$ \textbf{Sign}\,---\,ERK leads strain;\par\vskip0.1cm
      $\bullet$ \textbf{Lag}\,---\,$\sim$20 min, field-independent.
    \column{0.44\textwidth}
      \centering
      \begin{tikzpicture}[scale=0.9]
        \draw[->] (0,0) -- (4.2,0) node[right,font=\scriptsize]{lag (min)};
        \draw[->] (0,-0.2) -- (0,2.2) node[above,font=\scriptsize]{corr.};
        \draw[accentcolor, thick, domain=0:4, samples=60]
          plot (\x, {1.8*exp(-(\x-2)*(\x-2)/0.6)});
        \draw[dashed] (2,0) -- (2,1.8);
        \node[font=\scriptsize] at (2,-0.3) {20};
      \end{tikzpicture}
      \figcap{Fig 1\;cross-correlation peak}
  \end{columns}
\end{frame}

% ===== 3 · Mechanism =====
\sectiondivider{3}{Mechanism}

\begin{frame}
  \frametitle{One feedback term closes the loop}
  \vskip0.1cm
  Couple ERK activity $E$ to local strain $s$ with a single mechanical feedback $f(s)$:
  \vskip0.1cm
  \[ \dot E = -\,kE + f(s), \qquad \dot s = \alpha E - \beta s. \]
  \vskip0.1cm
  For a saturating $f$, the fixed point loses stability through a Hopf bifurcation,
  and coupling to neighbors turns the oscillation into a \alert{traveling wave}.
\end{frame}

\begin{frame}
  \frametitle{The minimal model, in three numbers}
  \vskip0.2cm
  \statrow{2}{coupled variables}{1}{feedback term}{0}{fitted parameters}
  \vskip0.6cm
  \begin{center}\small
  Everything below follows from these choices alone.
  \end{center}
\end{frame}

% ===== 4 · Prediction & test =====
\sectiondivider{4}{Prediction \& test}

\begin{frame}
  \frametitle{Three predictions we can falsify}
  \vskip0.1cm
  \hyporow{H1}{ERK leads strain by a fixed lag, set by $k$.}
          {H2}{Wave speed grows with feedback gain $\alpha$.}
          {H3}{Blocking feedback halts propagation.}
\end{frame}

\begin{frame}
  \frametitle{Model and experiment agree without tuning}
  \vskip0.1cm
  \begin{center}\small
  \begin{tabular}{@{}lcc@{}}
    \toprule
    & \textbf{Experiment} & \textbf{Model} \\ \midrule
    Phase lag (min)              & $19 \pm 3$ & $21$ \\
    Wave speed ($\mu$m/min)      & $2.4$      & $2.2$ \\
    Feedback block $\Rightarrow$ arrest & yes & yes \\
    \bottomrule
  \end{tabular}
  \end{center}
  \vskip0.2cm
  Both numbers fall within experimental error, with no free parameters.
\end{frame}

\statementframe{A single mechanical feedback is enough to sustain ERK waves ---
and it predicts the measured phase lead.}

\thanksframe{Thank you}{jane.roe@example.edu \quad$\cdot$\quad Biophysics Seminar}

\end{document}
